\documentclass[12pt]{article}
\usepackage[utf8]{inputenc}
\usepackage[T1]{fontenc}
\usepackage{hyperref}
\title{Waray Language Resources for K-12 Curriculum Development}
\author{Compiled by Claude Sonnet 5.5 (claude-sonnet-5-5) for Andr\'as Kornai}
\date{}
\begin{document}
\maketitle

\section{Introduction}
This dossier concerns Waray (ISO 639-3 wrz, Glottocode wara1290), an Australian Aboriginal language of the Gunwinyguan family. It is also spelled Warray. It is not the Philippine language Waray-Waray. It was spoken near the Adelaide River in the Northern Territory. Sources disagree on when it died out: Wikipedia gives about 2000 and ASJP gives 1990. The Endangered Languages catalogue lists it as critically endangered, which may reflect an older entry. No native speakers are known to survive, so any curriculum would have to be a revitalization effort built from archival documentation.

\section{Primary Educational Resources}
Searches found no primers, textbooks or school programmes for Waray. Revitalization teaching materials would have to be developed from the documentation listed below.

\section{Linguistic and Reference Materials}
\begin{itemize}
\item Wikipedia, Warray language: overview of location, classification and extinction date. \url{https://en.wikipedia.org/wiki/Warray_language}
\item AIATSIS Austlang entry N25 (Warray): language record from the Australian national catalogue. \url{https://collection.aiatsis.gov.au/austlang/language/N25}
\item Wikipedia, Gunwinyguan languages: family context. \url{https://en.wikipedia.org/wiki/Gunwinyguan_languages}
\item ASJP database entry for Waray. \url{https://asjp.clld.org/languages/WARAY}
\item Endangered Languages catalogue entry for Warray, with a bibliography. \url{https://www.endangeredlanguages.com/lang/2025?hl=es}
\item The Oxford Guide to Australian Languages, a chapter in the search results that places Warray in its family. I did not open it. \url{https://academic.oup.com/book/46706/chapter/410240832}
\item Mark Harvey (1986), \emph{Ngoni Waray amungal-yang: the Waray language from Adelaide River}, is cited as the key documentation. No URL for it was seen in the search results.
\end{itemize}

\section{Religious and Community Resources}
No scripture translations or community organisations specific to Waray were found.

\section{Digital and Audio Resources}
No audio, video or software resources for Waray were found.

\section{Curriculum Development Context}
Very little is publicly documented. A heritage curriculum would have to start from Harvey's 1986 work and the AIATSIS holdings, and ideally involve descendant communities and Aboriginal language centres in the Northern Territory. The searches did not identify such a centre, so I name none. Curriculum materials should be modest in scope, such as names, words and place names, until more documentation is located.
\end{document}